\documentclass[10pt,a4paper]{amsart}
\usepackage[margin=19mm]{geometry}
\usepackage{amsmath,amssymb,mathtools}
\usepackage[scr=rsfs,cal=euler]{mathalfa}
\usepackage{xfrac}
\usepackage[unicode,hidelinks,bookmarks=false]{hyperref}
\newcommand{\ie}{i.e.\ }
\newcommand{\cf}{cf.\ }
\newcommand{\resp}{respectively\ }
\newcommand{\Subsection}{\subsection}
\providecommand{\nobreakdash}{\nobreak\hbox{-}}
\setlength{\emergencystretch}{2em}
\multlinegap0pt
\usepackage{amssymb}
\newtheorem{theorem}{Theorem}
\title{The Equations of General Hassett maximal Cubic Fourfolds}
\author{Elad Gal, Howard Nuer}
\date{Source: arXiv:2605.09531v1 (CC BY 4.0)}
\begin{document}
\maketitle

\noindent\emph{Source and licence.} The Equations of General Hassett maximal Cubic Fourfolds. Version arXiv:2605.09531v1 (CC BY 4.0), distributed by arXiv under the Creative Commons Attribution 4.0 International licence, http://creativecommons.org/licenses/by/4.0/, as recorded in the arXiv licence metadata for this exact version and shown on the abstract page https://arxiv.org/abs/2605.09531v1. Full licence notice: \texttt{licenses/arxiv-2605.09531-CC-BY-4.0.txt}.\par\smallskip

\noindent\textit{Editorial scope.} Counted unit: the article's theorem theorem (source0.tex lines 103--106). The article's complete original proof of it is delivered with it (source0.tex lines 108--122, one \texttt{proof} environment of 15 lines). No mathematics is written, repaired or re-derived: the statement and the proof are the article's own lines, in the article's own notation, and no retained line is substituted or re-flowed.  No further source-local context is needed: every cross-reference of every retained line resolves inside the two delivered spans.  Nothing was omitted from any retained span, so the package carries no omission record.  No retained line contains a citation, so the wrapper carries no bibliography and no bibliography entry.  No retained line contains a figure environment, an image-inclusion command or drawing code, so no image is delivered and the package declares no asset dependency.  Editorial formatting only, in addition: an AMS \texttt{amsart} preamble that restates the article's own declarations and macros used by the retained lines, namely the environments \texttt{theorem} on separate counters, together with the standard packages those lines need.  The retained environments print in this document as Theorem 1 in order of appearance, whereas the article prints the same items with the numbering of its own full document.  The complete native source of the article as distributed in the arXiv e-print for this exact version is bundled unchanged as \texttt{sources/200230/source0.tex}.
\begin{theorem}
Let $F$ be a smooth homogeneous cubic polynomial defining a fourfold $V_+(F) \in \mathcal{N}$ in $\mathbb{P}^5 = \mathbb{P}\mathbb{C}^6$. Then there exists a basis of linear forms $\{x,y,z,u,v,w\}$ for the dual space $(\mathbb{C}^6)^*$ and constants $a,b \in \mathbb{C}$ such that:
$$F \in ( x,y,z ) \cap ( x,y,u ) \cap ( x,z,v ) \cap ( v-by,u-az,w )$$
\end{theorem}

\begin{proof}
Let $X = V_+(F) \in \mathcal{N}$, and let $P_1,P_2,P_3,P_4 \subset \mathbb{P}^5$ be the four planes exhibiting the required incidence profile. Each plane $P_i$ is the projectivization of a 3-dimensional vector subspace $V_i \subset \mathbb{C}^6$. The defining ideal $I(P_i)$ is generated by the 3-dimensional subspace of linear forms in $(\mathbb{C}^6)^*$ that vanish on $V_i$.

Since $P_1$ is a plane, we can choose three linearly independent linear forms $x,y,z \in (\mathbb{C}^6)^*$ such that $I(P_1) = ( x,y,z )$. 

Because the intersection $P_1 \cap P_2$ is a line, the sum $I(P_1) + I(P_2)$ must be generated by exactly 4 independent linear forms. Therefore, $I(P_2)$ shares a 2-dimensional subspace of linear forms with $I(P_1)$. We can choose the generators $x,y$ to span this intersection and extend with a new independent linear form $u$ such that $I(P_2) = ( x,y,u )$.

Similarly, $P_1 \cap P_3$ is a line, so $I(P_3)$ also shares a 2-dimensional subspace of linear forms with $I(P_1)$. For the intersection $P_2 \cap P_3$ to be exactly a point, the sum $I(P_2) + I(P_3)$ must be generated by 5 independent linear forms. This forces the shared forms between $I(P_1)$ and $I(P_3)$ to be spanned by $x$ and a form independent of $y$, which we take to be $z$. Extending with a fifth independent linear form $v$, we obtain $I(P_3) = ( x,z,v )$. Notice that $I(P_2) + I(P_3) = ( x,y,z,u,v )$, which correctly cuts out a point.

Finally, because $P_1 \cap P_4 = \emptyset$, the sum $I(P_1) + I(P_4)$ must be generated by the entire 6-dimensional dual space $(\mathbb{C}^6)^*$. This allows us to choose the remaining basis element $w$ such that $\{x,y,z,u,v,w\}$ forms a complete basis for $(\mathbb{C}^6)^*$. By applying suitable linear changes of basis to the generators of $I(P_4)$ while preserving our established basis elements, we can put the ideal of $P_4$ into the canonical form:
$$I(P_4) = ( v-by,u-az,w )$$
for some constants $a,b \in \mathbb{C}$. Because the fourfold $X = V_+(F)$ contains all four planes $P_i$, the polynomial $F$ must vanish on each $P_i$, placing it precisely in the intersection of their respective ideals:
$$F \in ( x,y,z ) \cap ( x,y,u ) \cap ( x,z,v ) \cap ( v-by,u-az,w ).$$
This completes the proof.
\end{proof}

\end{document}
